\section{Verification and trust boundary}\label{sec:verification}

The principal results are conceptual proofs.  This paper contains no
paper-facing finite census, generated certificate, or numerical experiment.
Its trust boundary has three layers: complete manuscript proofs, selected
kernel-checked algebraic cores, and standard results imported by citation.

\begin{table}[H]
\centering
\footnotesize
\caption{Trust boundary for the theorem families retained in this paper.
``Core formalization'' means that the named algebraic step is kernel checked;
it does not promote the complete manuscript theorem to a formal theorem.}
\label{tab:trust-boundary}
\begin{tabular}{@{}L{0.27\linewidth}L{0.34\linewidth}L{0.31\linewidth}@{}}
\toprule
Result family & Manuscript proof & Formal status \\
\midrule
AME support counting and complete Weyl-basis rigidity &
Sections~\ref{sec:dictionary}--\ref{sec:rigidity} &
coordinate squeeze, unique local labels, and the abstract axis criterion are
kernel checked; the concrete additive-state marginal composition is not \\
Minimum-support transition maps and phase correction &
Section~\ref{sec:atlas} &
minimum-support generation, the abstract holonomy centralizer, and the
stabilizer-character algebra are kernel checked; the state-ray correction is
a manuscript step; the systematic half-set form, block-minor criterion,
optimal direct sum, and finite recognition bounds are manuscript only \\
Marginal certification and stochastic conversion &
Section~\ref{sec:atlas} & manuscript only \\
Local endomorphism algebra, five prime-field types, and low-party theorem &
Section~\ref{sec:endomorphism-algebra} and
Appendix~\ref{app:endomorphism-code-structures} &
manuscript only; the existing abstract holonomy-centralizer core does not
formalize the intrinsic algebra identification, five-type composition,
low-party commutativity, or module consequences \\
Encoder Clifford no-go & Section~\ref{sec:atlas} &
the Choi orientation and Clifford closure are kernel checked; the complete
arbitrary-additive encoder theorem is assembled in the manuscript \\
Cleaning, Weyl--Fourier rounding, branch selection, and residual stability &
Section~\ref{sec:quantitative} &
manuscript only; neither the leakage constant \(8\), the Fourier rounding
lemma, the balanced-cut estimate, nor the global radius is formalized \\
Relative exact-base intertwiner reduction &
Section~\ref{sec:quantitative} &
exact line transport, defect invariance, product-intertwiner composition, and
lossless transfer of a supplied threshold and coefficient are kernel checked;
the cleaning inputs remain manuscript proofs \\
Robust transition compatibility and character obstruction &
Section~\ref{sec:quantitative} & manuscript only \\
Two-uniform discreteness and local stability & Appendix B &
only the generator splitting, single-exponential identity, polarized
second-moment identity, and absence of a nonscalar continuous symmetry are
kernel checked; the local quadratic constants are manuscript only \\
\bottomrule
\end{tabular}
\end{table}

Table~\ref{tab:trust-boundary} is a disclosure of partial formalization,
not an independently reproducible formal artifact for the principal theorems.
The selected cores were developed in the Lean namespace
\texttt{RelativeConicArcs.AMELU}.  The paper does not supply a self-contained
replay contract matching its statements to a pinned public set of formal
declarations, and claims no end-to-end formalization of any theorem family.
All principal results are proved in the manuscript independently of these
partial formal developments.

The cleaning theorem uses three cited exact inputs: stabilizer cleaning,
localization of nested commutators, and the exact three-cleanable-region
Clifford obstruction.  The manuscript proves the leakage estimate, the
finite Weyl--Fourier concentration step, the stabilizer overlap gap, and
their AME composition.  Its closed-form deductions do not depend on a
computational search or floating-point evaluation.

The quantitative appendix serves as a mechanism comparison.  None of its
results is used to prove Theorem~\ref{thm:quantitative-rounding}.
